\section{Related Work}
\label{sec:related}

\textbf{Claim extraction and factual verification.} Claimify~\cite{metropolitansky2025claimify} extracts atomic claims from LLM text with 99\% entailment. FActScore~\cite{min2023factscore} measures factual precision against Wikipedia. VeriFact~\cite{wei2024verifact} verifies clinical claims against health records. CCV extends this line of work to source code, with the key difference that verification is fully deterministic.

\textbf{LLM hallucination detection.} SelfCheckGPT~\cite{jiang2023selfcheck} detects hallucination via multi-sample consistency. Survey works~\cite{ji2023hallucination, zhang2023siren, liu2023trustworthy} categorize hallucination types and mitigation strategies. CCV operates downstream of hallucination detection: rather than detecting whether the LLM \emph{might} be hallucinating, CCV checks whether specific claims \emph{are} correct.

\textbf{LLM for code.} Codex~\cite{chen2021codex} demonstrated LLM code generation capabilities. Subsequent work studied security implications~\cite{pearce2022security} and automated repair~\cite{fan2023automated}. These works focus on generated code quality; CCV focuses on the accuracy of LLM \emph{reasoning about} existing code.

\textbf{Tool use verification.} Tool Receipts~\cite{wang2024toolreceipts} propose execution traces for verifiable tool use. CCV is complementary: it verifies the LLM's conclusions about code regardless of whether tools were used to reach those conclusions.

\textbf{Software engineering verification.} Static analysis tools (linters, type checkers) verify code properties but operate on code, not on natural language claims about code. CCV bridges this gap by translating LLM assertions into checkable code-level queries.
